% =====================================================================
% §1.4 — Lista de materiales (BOM) — Fase 1 (estructura física)
% =====================================================================
\section{Lista de materiales (BOM) --- Fase 1}
\label{sec:bom}

\subsection*{Resumen de la sección}
Solo se incluyen los componentes \textbf{físicos del rack y los
contenedores}. El BOM electrónico (sensores, ESP32, cámaras,
humidificadores) se consolida en la Fase 8 después de cerrar las
decisiones de las Fases 2 y 3, para evitar re-cotización.

Los precios están deliberadamente vacíos: cada renglón se marca con
\cotizar{descripción exacta a buscar} y se replica en
\texttt{cotizaciones\_pendientes.md} con la búsqueda sugerida por
proveedor (MercadoLibre MX, Home Depot, Steren, AliExpress).

% ---------------------------------------------------------------------
\subsection{Estructura del rack}
\label{sec:bom:rack}

% TODO[Fase 1 — relleno]:
% Tabla con columnas:
%   | Ítem | Descripción | Cant. | Económico (MXN) | Profesional (MXN) | Proveedor sugerido |
% Renglones esperados:
%   - Rack metálico \paramNiveles{} niveles, \paramHmax x \paramWmax x \paramDmax cm
%   - Anclaje a pared (opcional, antivuelco)
%   - Niveladores de patas (4 uds)
%   - Tapete antiderrame para nivel inferior
% Usar la macro \cotizar{...} en lugar de inventar precios.

\subsection{Contenedores y sellado}
\label{sec:bom:contenedores}

% TODO[Fase 1 — relleno]:
% Renglones esperados:
%   - Tote hermético \paramVolTote\,L con tapa (x\paramNcontenedores)
%   - Empaque de silicona alimenticia grado FDA (rollo)
%   - Prensaestopas PG7/PG9/PG11 según pasamuros
%   - Arandelas EPDM
%   - Tornillería inoxidable M3/M4 para fijación de sensores
%   - Tyvek o filtro microporoso para perforaciones (m^2)

\subsection{Herramientas específicas requeridas}
\label{sec:bom:herramientas}

% TODO[Fase 1 — relleno]:
%   - Sierra cilíndrica (hole saw) de diámetros derivados de §1.2.1
%   - Punta cónica para abrir pasamuros en plástico HDPE
%   - Termómetro infrarrojo para mapeo térmico del rack
%   - Higrómetro de referencia para calibración cruzada con SHT35
% Marcar las que ya tengo (laboratorio I+D) vs. las que faltan comprar.

\subsection{Política de placeholders de precio}
\label{sec:bom:politica}

\textbf{Regla estricta:} ningún precio se inventa.
\begin{itemize}[leftmargin=*]
  \item Si el componente fue cotizado por el operador, se sustituye el
        \cotizar{} por el valor real con fecha y URL.
  \item Mientras no haya cotización, el renglón permanece marcado y
        aparece replicado en \texttt{cotizaciones\_pendientes.md}.
  \item Las dos columnas (económico / profesional) se consolidan solo
        cuando ambas tienen al menos una cotización real cada una.
\end{itemize}

\subsection{Criterios de aceptación}
\label{sec:bom:aceptacion}

\begin{itemize}[leftmargin=*]
  \item Todo renglón sin precio aparece en
        \texttt{cotizaciones\_pendientes.md} con búsqueda sugerida.
  \item La cantidad de cada ítem está justificada (no se piden
        "varios", se piden números concretos derivados de \S\ref{sec:rack}
        y \S\ref{sec:contenedores}).
\end{itemize}
